\subsection{Milestone 3 -- Analisi dell’impatto dei code smell}

La Milestone 3 risponde alla domanda: \textit{quante classi buggy sarebbero potute essere
prevenute se non avessero avuto code smell?}

\subsubsection{Scelta del classificatore (BClassifier)}

La scelta di BClassifier è critica: un modello debole farebbe sì che la differenza
tra le predizioni su B+ e B rifletta il rumore del modello più che l'effetto reale
degli smell. Dall'analisi dei risultati di M2 (Sezione~3.2) emerge \textbf{RF + No FS + SMOTE}
come configurazione con il miglior bilanciamento complessivo su Kappa ($0{,}61$),
AUC ($0{,}94$) e NPofB20 ($0{,}88$).

\subsubsection{Costruzione dei dataset}

Dal dataset \textbf{A} (15.338 snapshot, 19 metriche, etichetta isBuggy) vengono derivati
tre dataset usando la metrica \texttt{Previous\_Release\_Code\_Smells} (NSmells):
\begin{itemize}[noitemsep]
  \item \textbf{B+}: istanze con NSmells $> 0$ (12.792 istanze, 83,4\% del totale).
  \item \textbf{C}: istanze con NSmells $= 0$ (2.546 istanze, 16,6\%).
  \item \textbf{B}: copia identica di B+ con NSmells forzato a 0.
\end{itemize}

\subsubsection{Training e predizione}

BClassifier è addestrato sull'intero A (\texttt{buildClassifier(A)}), con SMOTE applicato
internamente da \texttt{FilteredClassifier} solo sul training. BClassifierA predice A, B+,
B e C; il confronto istanza per istanza tra B+ e B isola l'effetto dell'azzeramento di NSmells.

Questa configurazione introduce un \textit{data leakage parziale}: B+ è costruito a partire
dalle stesse istanze di A con NSmells azzerato, e il modello è addestrato su A intero.
Le predizioni su B+ hanno sono quindi ottimistiche. Il leakage
è tuttavia accettabile per lo scopo dell'analisi: l'obiettivo non è stimare la capacità
predittiva del classificatore su dati nuovi, ma quantificare \textit{quanto cambiano le
predizioni} quando si rimuovono i code smell a parità di tutte le altre feature.

\subsubsection{Metriche di prevenzione}

Un'istanza è \textit{prevenuta} se BClassifierA la predice BUGGY su B+ e CLEAN su B.
Tre metriche quantificano la prevenzione:
\begin{itemize}[noitemsep]
  \item \textbf{Prevented\_Total}: numero assoluto di istanze prevenute.
  \item \textbf{Prevented\_Proportion}: Prevented\_Total / Actual\_Buggy\_in\_A.
  \item \textbf{Prevented\_Of\_Predictable}: Prevented\_Total / Predicted\_Buggy\_in\_B+.
\end{itemize}

\subsubsection{Analisi di ridondanza: correlazione e ablation study (M3.1)}

Due analisi supplementari spiegano empiricamente \textit{perché} l'impatto di NSmells
è basso.

\paragraph{Correlazione di Spearman.}
Per ogni feature (escluse NSmells e isBuggy) viene calcolata la correlazione $\rho$ di
Spearman con NSmells. Questa è una correlazione
\textit{feature a feature}, distinta dalla FS di M2 (feature a label): l'obiettivo non è
identificare feature poco predittive, ma quantificare la ridondanza informativa di NSmells.

\paragraph{Ablation study.}
10$\times$10-fold CV con RF + No FS + SMOTE su A a 19 feature e su A a 18 feature (senza
NSmells). Il $\Delta$ misura il contributo \textit{marginale} di NSmells.
